%% ================================================================
%%  中文论文 LaTeX 模板（XeLaTeX + ctex）
%%  适合：国内期刊论文 / 学位论文 / 研究报告
%%  编译：xelatex main.tex   （跑两遍）
%% ================================================================
\documentclass[UTF8,a4paper,11pt]{ctexart}

\usepackage{geometry}
\geometry{left=2.5cm,right=2.5cm,top=2.5cm,bottom=2.5cm}

\usepackage{amsmath,amssymb}
\usepackage{booktabs}
\usepackage{array}
\usepackage{tabularx}
\usepackage{ragged2e}
\usepackage{enumitem}
\usepackage{xcolor}
\usepackage{titlesec}
\usepackage{setspace}
\usepackage{fancyhdr}
\usepackage{hyperref}
\hypersetup{colorlinks=true,linkcolor=black,urlcolor=blue!50!black,citecolor=black}

%% 行距
\linespread{1.45}
%% 首行缩进两字
\setlength{\parindent}{2em}
\setlength{\parskip}{0.3em}

%% ⭐ 中文章节编号：一二三 / （一）（二）
%%    注意：不能用 \ctexset{section/number=...}，会被 titlesec 覆盖
\titleformat{\section}{\Large\heiti}{\chinese{section}、}{0.3em}{}
\titleformat{\subsection}{\large\heiti}{（\chinese{subsection}）}{0.2em}{}
\titlespacing{\section}{0pt}{1.5em}{0.8em}
\titlespacing{\subsection}{0pt}{1em}{0.5em}

%% 页脚页码
\pagestyle{fancy}
\fancyhf{}
\renewcommand{\headrulewidth}{0pt}
\fancyfoot[C]{\small\thepage}

%% 摘要环境
\newenvironment{cnabstract}{%
  \begin{center}\small\heiti 摘\quad 要\end{center}
  \begin{quote}\small\setstretch{1.25}%
}{%
  \end{quote}
}

%% 关键词
\newcommand{\keywords}[1]{%
  \par\vspace{0.6em}\noindent{\small\heiti 关键词：}#1\par\vspace{1.2em}%
}

%% 强调（中文论文里加粗即可，不用斜体）
\newcommand{\kw}[1]{\textbf{#1}}


\begin{document}

%% ---------------- 标题区 ----------------
\begin{center}
  {\LARGE\heiti 论文标题写这里}\\[0.5em]
  {\large\heiti ——副标题写这里}\\[1.2em]
  {\normalsize 作者姓名}\\[0.4em]
  {\small （作者单位）}\\[0.3em]
  {\small 2026 年 11 月}
\end{center}

\vspace{1em}

%% ---------------- 摘要 ----------------
\begin{cnabstract}
摘要正文写这里。建议三到五段：第一段提出问题，第二段说明方法，
第三段给出结论，第四段说明贡献与限定。

注意中文论文的摘要通常是**报道性摘要**，要说清"做了什么、怎么做的、
得到什么结果"，而不是只介绍背景。
\end{cnabstract}

\keywords{关键词一；关键词二；关键词三；关键词四}

%% ---------------- 正文 ----------------
\section{引言}

正文从这里开始。首行自动缩进两字。

这是一个段落示例。中文排版需要注意标点挤压、行首行尾禁则、以及中西文混排时的间距，
xeCJK/ctex 已经处理好了，不需要手动加空格。

\subsection{研究背景}

小节用「（一）（二）」编号。

\subsection{本文贡献}

\begin{enumerate}[leftmargin=2em,itemsep=0.3em]
  \item 贡献一；
  \item 贡献二；
  \item 贡献三。
\end{enumerate}

\section{方法}

\subsection{公式}

行内公式：$E = mc^2$。

独立公式：
\begin{equation}
\cos\theta = \frac{\mathbf{a} \cdot \mathbf{b}}{\lVert \mathbf{a}\rVert \, \lVert \mathbf{b}\rVert}
\end{equation}

\subsection{表格}

中文论文习惯用三线表：

\begin{center}
\begin{tabular}{@{}p{3.2cm}p{4.6cm}p{4.6cm}@{}}
\toprule
 & 对照组 & 实验组 \\
\midrule
指标一 & 数值 & 数值 \\
指标二 & 数值 & 数值 \\
指标三 & 数值 & 数值 \\
\bottomrule
\end{tabular}
\end{center}

\subsection{列表}

\begin{itemize}[leftmargin=2em,itemsep=0.2em]
  \item 要点一；
  \item 要点二。
\end{itemize}

\section{结论}

结论写这里。不要简单复述摘要，要指出：
\begin{enumerate}[leftmargin=2em,itemsep=0.3em]
  \item 得出了什么；
  \item 局限在哪；
  \item 下一步是什么。
\end{enumerate}

%% ---------------- 参考文献 ----------------
\vspace{1.5em}
\section*{参考文献}

\begin{list}{}{\leftmargin=2em\itemindent=-2em\itemsep=0.35em}
\small

\item 作者. 文献标题[J]. 期刊名, 年份, 卷(期): 起止页码.

\item 作者. 书名[M]. 出版地: 出版社, 年份.

\item Author A, Author B. Paper Title[C]//\textit{Proceedings of Conference}. Year: pages.

\end{list}

\end{document}
